\section{Revision notes}\label{app:revisions}

This appendix lists the changes from the first version of the paper. The main text states the
results as they now stand and does not refer to the earlier version.

\begin{proposition}[the version-1 capacity measure is vacuous]\label{prop:e2-v1-vacuous}
Version~1 measured an audit tower by a record universe with $|U(z)| = \capc(z)$ and, for every
occurring level $\ell$, a record set $R(\ell) \subseteq U(z)$ containing exactly $\ell$'s counted
records, with $R(\ell) \cap R(\ell') = \varnothing$ whenever $\ell$ and $\ell'$ occur and have
different indices. No E-system with an occurring level admits such a measure.
\end{proposition}

\begin{proof}
Let $\ell$ occur with index $k$, and let $\ell'$ be $\ell$ with index $k + 1$ and an empty target
list. Then $\ell'$ occurs, and the tag of the common own record lies in $R(\ell)$ and in $R(\ell')$,
contradicting disjointness.
\end{proof}

\begin{proposition}[the version-1 forcing certificate is uninhabited]\label{prop:e1-v1-vacuous}
Version~1 derived infinitely many strict self-extensions from a finite-forcing certificate that
asserts the conclusion for every repair family, every novelty predicate and every strictness
predicate whenever the challenge is new at every time. No system with inhabited state types admits
such a certificate.
\end{proposition}

\begin{proof}
Take the novelty predicate that always holds and the empty family. The certificate asserts that the
empty family has an entry beyond every bound, and it has none.
\end{proof}

\begin{enumerate}
  \item \emph{E1.} By \Cref{prop:e1-v1-vacuous}, the strict-self-extension statement of version~1
    had no instance. The actual-family certificate of \Cref{def:e1-forcing} replaces it, with the
    clocked instance of \Cref{thm:e1-clocked}; E1 is now a law.
  \item \emph{E2.} By \Cref{prop:e2-v1-vacuous}, the capacity bound of version~1 held only for
    systems without an occurring instrument level. The same measure was a hypothesis of forced
    stratification, of the level classification, of the regress stop of E3 and of the meta-audit bound
    of E8.1; none of these applies to a system with an occurring level. The declared measure of
    \Cref{def:declared-measure} replaces it, with the instance of \Cref{thm:e2-tower}. The version-1
    declarations remain in the library.
  \item \emph{E3.} The reinstatement condition of version~1 names a challenge history but does not
    use it, so it holds with an empty history. The history-linked form of \Cref{def:episode-linked}
    replaces it, and the device now executes its repair along its declared trajectory.
  \item \emph{E6.} Version~1 called the witnessed allocation ``certified optimal''. The multiplier
    witness yields the consequences of \Cref{thm:e6-kkt}, not optimality; optimality holds under the
    hypotheses of \Cref{prop:e6-sufficient}, and is checked in Lean for capped linear objectives
    (\Cref{thm:e6-caps}).
  \item \emph{E12 and E13.} In version~1 both were templates whose certificates had budget,
    viability and completeness fields that no concrete model could establish. They are now laws with
    the semantic definitions of \Cref{def:e12-setting,def:e13-transport}; the version-1 Templates are
    kept as the full classifications.
  \item \emph{E15.} The budget-reallocation identity was a field of the allocation record in
    version~1; it is now derived from the exhaustion equations.
  \item \emph{Presentation.} The paper leads with the laws and their instances, gives the
    certificate specifications afterwards, and moves the concordance, field lists and registers to
    the supplement.
\end{enumerate}
